\section{结论}
现有工作在设计和改进鸟瞰视图（BEV）识别模型的检测器方面付出了大量努力，但它们通常局限于特定的预训练骨干网络，未进一步探索。
本文旨在释放现代图像骨干网络（Image Backbone）在BEV模型上的全部潜力。
本文将通用2D图像骨干网络适配的困难归因于BEV检测器的优化问题。
为解决此问题，本文通过从额外的透视3D检测头添加辅助损失（Auxiliary Loss），将透视监督（Perspective Supervision）引入BEV模型。
此外，本文将两个检测头整合为两阶段（Two-stage）检测器，即 BEVFormer v2。
功能完整的透视头提供第一阶段物体提议（Object Proposal），编码为BEV头的物体查询（Object Query）用于第二阶段预测。
大量实验验证了本文方法的有效性和泛化性。
透视监督引导2D图像骨干网络感知自动驾驶的3D场景，帮助BEV模型实现更快的收敛和更好的性能，并且适用于广泛的骨干网络。
此外，本文成功将大规模骨干网络适配到 BEVFormer v2，在 nuScenes 数据集上取得了新的SOTA结果。
本文的工作为未来研究者探索更好的BEV模型图像骨干网络设计铺平了道路。

\noindent\textbf{局限性.} 由于计算和时间限制，本文目前未在更大规模的图像骨干网络上测试方法。本文已在多种骨干网络上完成了初步验证，未来将扩展到更大的模型规模。
